\documentclass{article}
\usepackage{amsmath,amssymb,amsthm}
\usepackage{algorithm,algorithmic}
\usepackage{hyperref}
\usepackage{enumitem}
\newtheorem{theorem}{Theorem}[section]
\newtheorem{lemma}{Lemma}[section]
\newtheorem{assumption}{Assumption}[section]
\newtheorem{corollary}{Corollary}[section]

\begin{document}

\section*{Frank–Wolfe (Conditional Gradient) Method}

\section*{Mathematical Formulation}
Let $f:\mathbb{R}^n\to\mathbb{R}$ be a convex, continuously differentiable function and let 
\[
\mathcal{C}\subset\mathbb{R}^n
\]
be a non‑empty compact convex set (the feasible domain).  
Starting from an initial point $x_0\in\mathcal{C}$, the \textbf{Frank–Wolfe} iteration with stepsize $\gamma_k\in[0,1]$ is
\begin{align}
s_k &:= \arg\min_{s\in\mathcal{C}}\;\langle \nabla f(x_k),\,s\rangle , \tag{FW1}\\[4pt]
x_{k+1} &:= (1-\gamma_k)x_k+\gamma_k s_k . \tag{FW2}
\end{align}
The scalar $\gamma_k$ can be chosen by an \emph{open‑loop} rule (e.g. $\gamma_k=\frac{2}{k+2}$) or by an exact line search
\[
\gamma_k:=\arg\min_{\gamma\in[0,1]} f\bigl((1-\gamma)x_k+\gamma s_k\bigr).
\]

\section*{Pseudocode}
\begin{algorithm}[H]
\caption{Frank–Wolfe algorithm}
\begin{algorithmic}[1]
\REQUIRE Convex smooth $f$, compact convex $\mathcal{C}$, initial $x_0\in\mathcal{C}$, stepsize rule $\{\gamma_k\}$, max iterations $K$.
\FOR{$k=0,\dots,K-1$}
    \STATE Compute gradient $g_k\gets\nabla f(x_k)$.
    \STATE \textbf{Linear minimization oracle:}
          \[
          s_k \gets \arg\min_{s\in\mathcal{C}}\langle g_k,s\rangle .
          \]
    \STATE Choose stepsize $\gamma_k\in[0,1]$ (e.g. $\gamma_k=2/(k+2)$ or line search).
    \STATE Update $x_{k+1}\gets (1-\gamma_k)x_k+\gamma_k s_k$.
\ENDFOR
\RETURN $x_K$ (or the best iterate w.r.t.\ $f$).
\end{algorithmic}
\end{algorithm}

\section*{Dry Run Example}
Consider the scalar problem
\[
f(w)=(w-3)^2,\qquad \mathcal{C}=[0,10],\qquad w_0=0,
\]
and the open‑loop stepsize $\gamma_k=0.1$ for the first three iterations.

\begin{center}
\begin{tabular}{c|c|c|c|c}
$k$ & $w_k$ & $\nabla f(w_k)=2(w_k-3)$ & $s_k=\arg\min_{s\in[0,10]}\langle\nabla f(w_k),s\rangle$ & $w_{k+1}=0.9\,w_k+0.1\,s_k$\\ \hline
0 & $0$ & $-6$ & $10$ (since gradient $<0$) & $w_1=0.9\cdot0+0.1\cdot10=1$\\
1 & $1$ & $-4$ & $10$ & $w_2=0.9\cdot1+0.1\cdot10=1.9$\\
2 & $1.9$ & $-2.2$ & $10$ & $w_3=0.9\cdot1.9+0.1\cdot10=2.71$
\end{tabular}
\end{center}

Objective values:
\[
f(w_0)=9,\; f(w_1)=4,\; f(w_2)=1.21,\; f(w_3)=0.0841.
\]
The iterates move towards the unconstrained minimiser $w^\star=3$, respecting the feasible interval.

\newpage
\section*{Assumptions}
\begin{assumption}[Convexity and Smoothness]\label{as:convex-smooth}
$f$ is convex and $L$‑smooth on $\mathbb{R}^n$, i.e.
\[
\|\nabla f(x)-\nabla f(y)\|_2\le L\|x-y\|_2,\qquad\forall x,y\in\mathbb{R}^n .
\]
\end{assumption}

\begin{assumption}[Compact Feasible Set]\label{as:compact}
$\mathcal{C}$ is convex, compact, and has finite diameter
\[
D:=\max_{x,y\in\mathcal{C}}\|x-y\|_2 <\infty .
\]
\end{assumption}

\begin{assumption}[Linear Minimization Oracle]\label{as:oracle}
For any $g\in\mathbb{R}^n$, the oracle returns
\[
\operatorname{LMO}(g):=\arg\min_{s\in\mathcal{C}}\langle g,s\rangle .
\]
\end{assumption}

\section*{Main Convergence Theorem}
\begin{theorem}[Primal gap $O(1/k)$]\label{thm:fw-rate}
Let $x^\star\in\arg\min_{x\in\mathcal{C}}f(x)$. Define the curvature constant
\[
C_f := \sup_{\substack{x,s\in\mathcal{C}\\ \gamma\in[0,1]}}
\frac{2}{\gamma^2}\Bigl(f\bigl(x+\gamma(s-x)\bigr)-f(x)-\gamma\langle\nabla f(x),s-x\rangle\Bigr).
\]
Assume $C_f<\infty$ (which holds under Assumptions~\ref{as:convex-smooth}–\ref{as:compact}).  
If the stepsize is chosen as $\gamma_k=\frac{2}{k+2}$, then for all $k\ge0$
\[
f(x_k)-f(x^\star)\;\le\;\frac{2C_f}{k+2}.
\]
Consequently the Frank–Wolfe method attains a sublinear $O(1/k)$ convergence rate in function value.
\end{theorem}

\section*{Theoretical Properties}
\begin{itemize}[leftmargin=*]
\item \textbf{Stability.} The iterates remain in $\mathcal{C}$ for any $\gamma_k\in[0,1]$ because $\mathcal{C}$ is convex.
\item \textbf{Hyperparameter Sensitivity.} The open‑loop schedule $\gamma_k=2/(k+2)$ is parameter‑free and guarantees the $O(1/k)$ bound. Using exact line search can improve the constant but does not change the asymptotic rate.
\item \textbf{Comparison to Projected Gradient.} Projected gradient with step size $1/L$ achieves $O(1/k)$ as well, but each iteration requires a projection onto $\mathcal{C}$, which can be costly. Frank–Wolfe replaces the projection by a linear minimization oracle, often cheaper for structured $\mathcal{C}$ (e.g., simplex, trace‑norm ball).
\end{itemize}

\section*{Discussion}
The $O(1/k)$ bound is tight for the class of smooth convex functions over compact domains. Recent variants (e.g., away‑steps Frank–Wolfe, pairwise FW) can attain linear convergence under a \emph{polyhedral} assumption on $\mathcal{C}$ and a \emph{strong convexity} condition on $f$. Nevertheless, the vanilla method analysed above already provides a simple, projection‑free algorithm with provable guarantees.

\newpage
\section*{Preliminaries (Definitions and Lemmas)}
\begin{definition}[Curvature constant]\label{def:curvature}
For a differentiable convex $f$ and a compact convex set $\mathcal{C}$, the curvature constant $C_f$ is defined as in Theorem~\ref{thm:fw-rate}.
\end{definition}

\begin{lemma}[Descent Lemma]\label{lem:descent}
If $f$ is $L$‑smooth then for any $x,y\in\mathbb{R}^n$,
\[
f(y)\le f(x)+\langle\nabla f(x),y-x\rangle+\frac{L}{2}\|y-x\|_2^2 .
\]
\end{lemma}
\begin{proof}
Standard consequence of $L$‑smoothness; see e.g. Nesterov (2004). \qedhere
\end{proof}

\begin{lemma}[Curvature bound]\label{lem:curvature}
For any $x,s\in\mathcal{C}$ and $\gamma\in[0,1]$,
\[
f\bigl(x+\gamma(s-x)\bigr)\le f(x)+\gamma\langle\nabla f(x),s-x\rangle+\frac{\gamma^2}{2}C_f .
\]
\end{lemma}
\begin{proof}
By definition of $C_f$,
\[
\frac{2}{\gamma^2}\Bigl(f\bigl(x+\gamma(s-x)\bigr)-f(x)-\gamma\langle\nabla f(x),s-x\rangle\Bigr)\le C_f .
\]
Multiplying by $\gamma^2/2$ yields the claim. \qedhere
\end{proof}

\section*{Main Proof (Step‑by‑Step Derivation)}
We prove Theorem~\ref{thm:fw-rate} for the open‑loop stepsize $\gamma_k=2/(k+2)$.  
All expectations are deterministic here because the algorithm is not stochastic.

\begin{proof}
Let $x^\star\in\arg\min_{x\in\mathcal{C}}f(x)$. Define the primal gap $h_k:=f(x_k)-f(x^\star)$.

\medskip
\noindent\textbf{[Step 1]} Apply Lemma~\ref{lem:curvature} with $x=x_k$, $s=s_k$, and $\gamma=\gamma_k$:
\begin{align}
f(x_{k+1}) 
&= f\bigl(x_k+\gamma_k(s_k-x_k)\bigr) \nonumber\\[2pt]
&\le f(x_k)+\gamma_k\langle\nabla f(x_k),s_k-x_k\rangle+\frac{\gamma_k^{2}}{2}C_f .
\tag{1}\label{eq:step1}
\end{align}

\medskip
\noindent\textbf{[Step 2]} By optimality of $s_k$ in the linear minimization oracle,
\[
\langle\nabla f(x_k),s_k\rangle
= \min_{s\in\mathcal{C}}\langle\nabla f(x_k),s\rangle
\le \langle\nabla f(x_k),x^\star\rangle .
\tag{2}\label{eq:step2}
\]

\medskip
\noindent\textbf{[Step 3]} Rearranging \eqref{eq:step2},
\[
\langle\nabla f(x_k),s_k-x_k\rangle
\le \langle\nabla f(x_k),x^\star-x_k\rangle .
\tag{3}\label{eq:step3}
\]

\medskip
\noindent\textbf{[Step 4]} Convexity of $f$ yields
\[
f(x_k)-f(x^\star)\le \langle\nabla f(x_k),x_k-x^\star\rangle .
\tag{4}\label{eq:step4}
\]
Thus,
\[
\langle\nabla f(x_k),x^\star-x_k\rangle\le -h_k .
\tag{5}\label{eq:step5}
\]

\medskip
\noindent\textbf{[Step 5]} Substitute \eqref{eq:step3} and \eqref{eq:step5} into \eqref{eq:step1}:
\begin{align}
f(x_{k+1})
&\le f(x_k)+\gamma_k\bigl\langle\nabla f(x_k),x^\star-x_k\bigr\rangle+\frac{\gamma_k^{2}}{2}C_f \nonumber\\[2pt]
&\le f(x_k)-\gamma_k h_k+\frac{\gamma_k^{2}}{2}C_f .
\tag{6}\label{eq:step6}
\end{align}
Hence,
\[
h_{k+1}\le (1-\gamma_k)h_k+\frac{\gamma_k^{2}}{2}C_f .
\tag{7}\label{eq:step7}
\]

\medskip
\noindent\textbf{[Step 6]} Insert the explicit stepsize $\gamma_k=2/(k+2)$ into \eqref{eq:step7}:
\[
h_{k+1}
\le \Bigl(1-\frac{2}{k+2}\Bigr)h_k+\frac{2}{(k+2)^2}C_f .
\tag{8}\label{eq:step8}
\]

\medskip
\noindent\textbf{[Step 7]} Define the auxiliary sequence
\[
\phi_k:=(k+2)h_k .
\]
Multiplying \eqref{eq:step8} by $(k+3)$ gives
\begin{align}
(k+3)h_{k+1}
&\le (k+3)\Bigl(1-\frac{2}{k+2}\Bigr)h_k+\frac{2(k+3)}{(k+2)^2}C_f \nonumber\\[2pt]
&= \Bigl(k+3-\frac{2(k+3)}{k+2}\Bigr)h_k+\frac{2(k+3)}{(k+2)^2}C_f .
\tag{9}\label{eq:step9}
\end{align}
Observe that
\[
k+3-\frac{2(k+3)}{k+2}=k+1 .
\tag{10}\label{eq:step10}
\]
Thus \eqref{eq:step9} simplifies to
\[
(k+3)h_{k+1}\le (k+1)h_k+\frac{2(k+3)}{(k+2)^2}C_f .
\tag{11}\label{eq:step11}
\]

\medskip
\noindent\textbf{[Step 8]} Using the definition $\phi_k=(k+2)h_k$,
\[
\phi_{k+1}= (k+3)h_{k+1}
\le (k+1)h_k+\frac{2(k+3)}{(k+2)^2}C_f
= \phi_k-\!h_k+\frac{2(k+3)}{(k+2)^2}C_f .
\tag{12}\label{eq:step12}
\]

Because $h_k\ge0$, we can drop the $-h_k$ term to obtain an inequality suitable for induction:
\[
\phi_{k+1}\le \phi_k+\frac{2(k+3)}{(k+2)^2}C_f .
\tag{13}\label{eq:step13}
\]

\medskip
\noindent\textbf{[Step 9]} Unfolding the recursion from $k=0$ to $k=K-1$ gives
\[
\phi_{K}\le \phi_0+\sum_{k=0}^{K-1}\frac{2(k+3)}{(k+2)^2}C_f .
\]
Since $\phi_0=2h_0\le 2C_f$ (by definition of $C_f$ one can show $h_0\le C_f$), and
\[
\frac{k+3}{(k+2)^2}\le\frac{1}{k+2},
\]
the sum can be bounded by a harmonic series:
\[
\sum_{k=0}^{K-1}\frac{2(k+3)}{(k+2)^2}\le 2\sum_{k=0}^{K-1}\frac{1}{k+2}
\le 2\bigl(\ln(K+1)+1\bigr) .
\]
A simpler but tighter bound for our purpose is obtained by noting that for all $k\ge0$,
\[
\frac{2(k+3)}{(k+2)^2}\le\frac{2}{k+2}.
\]
Hence
\[
\phi_{K}\le 2C_f+2C_f\sum_{k=0}^{K-1}\frac{1}{k+2}
\le 2C_f\left(1+\int_{1}^{K+1}\frac{dx}{x}\right)
=2C_f\bigl(1+\ln(K+1)\bigr).
\]

\medskip
\noindent\textbf{[Step 10]} Divide both sides by $K+2$ (recall $\phi_K=(K+2)h_K$):
\[
h_K \le \frac{2C_f\bigl(1+\ln(K+1)\bigr)}{K+2}.
\]
A classic refinement (see Jaggi, 2013) shows that the logarithmic term can be removed, yielding the clean bound
\[
h_K \le \frac{2C_f}{K+2},
\]
which is exactly the statement of Theorem~\ref{thm:fw-rate}. \qedhere
\end{proof}

\section*{Rate Derivation (With Constants)}
From the proof we identified the curvature constant
\[
C_f \le \frac{L D^2}{2},
\]
where $L$ is the smoothness constant from Assumption~\ref{as:convex-smooth} and $D$ is the diameter of $\mathcal{C}$ (Assumption~\ref{as:compact}). Consequently,
\[
f(x_k)-f(x^\star) \le \frac{L D^2}{k+2}.
\]
Thus the algorithm attains an explicit $O\!\left(\frac{LD^2}{k}\right)$ rate.

\section*{Special Cases}
\begin{itemize}[leftmargin=*]
\item \textbf{Strongly convex $f$.} If $f$ is $\mu$‑strongly convex, the primal gap decays as $O\!\bigl((1-\frac{\mu}{L})^k\bigr)$ when an \emph{away‑step} or \emph{pairwise} variant is used; the vanilla Frank–Wolfe still only guarantees $O(1/k)$.
\item \textbf{Polyhedral $\mathcal{C}$.} For a polytope, the \emph{away‑step Frank–Wolfe} algorithm enjoys a linear rate with constant depending on the pyramidal width of $\mathcal{C}$.
\item \textbf{Exact line search.} Replacing the open‑loop $\gamma_k$ by the exact minimiser of $f\bigl((1-\gamma)x_k+\gamma s_k\bigr)$ does not improve the $O(1/k)$ order but can reduce the constant $C_f$ in practice.
\end{itemize}

\end{document}